% notes/artefacts/portfolio-construction-risk/portfolio-construction-models.tex
% Legacy file kept verbatim when it was migrated (notes/source-review-records.md, portfolio construction review). Its single map candidate,
% portfolio-construction, owns all its entries; the file sits beside transaction-cost models review's transaction-cost-models.tex under the
% section key, as a new file so that the earlier artefact stays unchanged. The compendium entries state each claim
% with its locator; the original wording lives here. Treat everything here as unchecked own working.
% Compile inside notes/artefacts/portfolio-construction-risk/ with  latexmk -pdf portfolio-construction-models.tex
\PassOptionsToPackage{table}{xcolor}
\documentclass[a4paper,oneside,openany]{book}
% Your own packages, macros and environments. The learning guide's entry environments,
% \citep and \qfsource are in house.tex, which compendium.tex loads after this file.

% -----------------------------
% Document Class & Encoding
% -----------------------------
\usepackage[utf8]{inputenc}  % Input encoding
\usepackage[T1]{fontenc}     % Font encoding

% -----------------------------
% Page Layout & Geometry
% -----------------------------
\usepackage[margin=2cm,foot=1cm]{geometry}

% -----------------------------
% Header & Footer Management
% -----------------------------
\usepackage{fancyhdr}
\pagestyle{fancy}

% -----------------------------
% Text Formatting & Spacing
% -----------------------------
\usepackage{ragged2e}
\usepackage[nodisplayskipstretch]{setspace}
\setstretch{1}  % Line spacing

% -----------------------------
% Listings for Code
% -----------------------------
\usepackage{listings}
\usepackage{xcolor}

% -----------------------------
% Tables, Figures & Graphics
% -----------------------------
\usepackage{tabularx}
\usepackage{multirow}
\usepackage{pdflscape}
\usepackage{graphicx}

% -----------------------------
% Algorithms
% -----------------------------
\usepackage{algorithm}       % For floating algorithm environment
\usepackage{algpseudocode}   % For pseudo-code inside algorithms

% -----------------------------
% Math Packages
% -----------------------------
\usepackage{amsmath,amssymb,amsthm,mathtools,bm}
\usepackage{mathrsfs}        % Script fonts
\usepackage{tikz-cd}         % Commutative diagrams
\usepackage{rotating}        % Rotated objects
\usepackage{float}           % Float placement
\usetikzlibrary{calc, arrows, positioning, arrows.meta}

% -----------------------------
% TikZ and its Libraries
% -----------------------------
\usepackage{tikz}
\usetikzlibrary{positioning} % Added to support "below=of" syntax
\usepackage{forest}          % For advanced tree diagrams
\usepackage{tikz-qtree}      % For vertical tree diagrams

% -----------------------------
% Miscellaneous Packages
% -----------------------------
\usepackage{xparse}          % For defining robust commands
\usepackage{enumitem}        % For customizing list layouts
\usepackage{makecell}        % For enhanced table cells
\usepackage{comment}		 % For commenting out sections

%%%%%%%%%%%%%%%%%%%%%%%%%%%%%%%%%%%%%%%%%%%%%%%%%%%%%%%%%%%%%%%%%%%%%%%%%%%%%%%
% Custom Math Commands and Shortcuts
%%%%%%%%%%%%%%%%%%%%%%%%%%%%%%%%%%%%%%%%%%%%%%%%%%%%%%%%%%%%%%%%%%%%%%%%%%%%%%%

% Matrix and Vector Formatting
\newcommand{\matr}[1]{\mathbf{#1}}      % For undergraduate algebra (alternative: \bm{#1})
\newcommand{\vect}[1]{\bm{#1}}

% Standard Number Sets
\newcommand{\N}{\mathbb{N}}
\newcommand{\R}{\mathbb{R}}
\newcommand{\Q}{\mathbb{Q}}
\newcommand{\C}{\mathbb{C}}
\newcommand{\Z}{\mathbb{Z}}

% Bold variables (customize as needed)
\newcommand{\x}{\mathbf{x}}
\newcommand{\F}{\mathbf{F}}
\newcommand{\f}{\mathbf{f}}
\newcommand{\y}{\mathbf{y}}
\renewcommand{\a}{\mathbf{a}}
\renewcommand{\b}{\mathbf{b}}
\renewcommand{\c}{\mathbf{c}}
\newcommand{\h}{\mathbf{h}}
\newcommand{\g}{\mathbf{g}}
\newcommand{\z}{\mathbf{z}}
\newcommand{\ze}{\mathbf{0}}

% Norms, Absolute Values, and Bracketing
\newcommand{\norm}[1]{\left\lVert#1\right\rVert}
\newcommand{\abs}[1]{\left\lvert#1\right\rvert}
\newcommand{\brk}[1]{\left[#1\right]}
\newcommand{\brc}[1]{\left\{#1\right\}}
\newcommand{\paren}[1]{\left(#1\right)}
\newcommand{\normop}[1]{\left\lVert#1\right\rVert_\text{op}}

% Miscellaneous Math Shortcuts
\newcommand{\LL}{\mathcal{L}}
\newcommand{\uni}{\overset{\text{uni}}{\to}}
\DeclareMathOperator{\diam}{diam}
\newcommand{\Prr}[1]{\text{Pr}\left(#1\right)}
\newcommand{\ceil}[1]{\lceil #1 \rceil}
\newcommand{\floor}[1]{\lfloor #1 \rfloor}
\newcommand\disteq{\stackrel{\textnormal{dist}}{=}}

% Evaluation command with options
\NewDocumentCommand{\evalat}{sO{\big}mm}{%
  \IfBooleanTF{#1}
    {\mleft. #3 \mright|_{#4}}
    {#3#2|_{#4}}%
}

%%%%%%%%%%%%%%%%%%%%%%%%%%%%%%%%%%%%%%%%%%%%%%%%%%%%%%%%%%%%%%%%%%%%%%%%%%%%%%%
% Color Definitions
%%%%%%%%%%%%%%%%%%%%%%%%%%%%%%%%%%%%%%%%%%%%%%%%%%%%%%%%%%%%%%%%%%%%%%%%%%%%%%%

\definecolor{dartmouthgreen}{rgb}{0.05, 0.5, 0.06}
\definecolor{egyptianblue}{rgb}{0.06, 0.2, 0.65}
\definecolor{dukeblue}{rgb}{0.0, 0.0, 0.61}
\definecolor{jazzberryjam}{rgb}{0.65, 0.04, 0.37}
\definecolor{magenta}{HTML}{EC008C}
\definecolor{darkmagenta}{rgb}{0.55, 0.0, 0.55}
\definecolor{deeppink}{rgb}{1.0, 0.08, 0.58}
\definecolor{codegreen}{rgb}{0,0.6,0}
\definecolor{codegray}{rgb}{0.5,0.5,0.5}
\definecolor{codepurple}{rgb}{0.58,0,0.82}
\definecolor{backcolour}{rgb}{0.95,0.95,0.92}

%%%%%%%%%%%%%%%%%%%%%%%%%%%%%%%%%%%%%%%%%%%%%%%%%%%%%%%%%%%%%%%%%%%%%%%%%%%%%%%
% Custom Shortcuts and Emphasis Commands
%%%%%%%%%%%%%%%%%%%%%%%%%%%%%%%%%%%%%%%%%%%%%%%%%%%%%%%%%%%%%%%%%%%%%%%%%%%%%%%

\newcommand{\isomorp}{\xrightarrow{\sim}}
\newcommand{\hlt}[1]{\textit{{\color{blue}#1}}}
\newcommand{\simpt}[1]{\textit{{\color{deeppink}#1}}}
\newcommand{\impt}[1]{\textit{{\color{red}#1}}}
\newcommand{\gcds}[1]{\textnormal{gcd}#1}
\newcommand{\mins}[1]{\textnormal{min}#1}
\newcommand{\maxs}[1]{\textnormal{max}#1}
\newcommand{\lcms}[1]{\textnormal{lcm}#1}
\newcommand{\degs}[1]{\textnormal{deg}#1}
\newcommand{\exps}[1]{\textnormal{exp}#1}
\newcommand{\tors}[1]{\textnormal{Tor}#1}
\newcommand{\homs}[1]{\textnormal{Hom}#1}
\newcommand{\anns}[1]{\textnormal{Ann}#1}
\newcommand{\defeq}{\vcentcolon=}
\DeclareMathOperator{\sign}{sign}

%%%%%%%%%%%%%%%%%%%%%%%%%%%%%%%%%%%%%%%%%%%%%%%%%%%%%%%%%%%%%%%%%%%%%%%%%%%%%%%
% Theorem and Definition Environments
%%%%%%%%%%%%%%%%%%%%%%%%%%%%%%%%%%%%%%%%%%%%%%%%%%%%%%%%%%%%%%%%%%%%%%%%%%%%%%%
\theoremstyle{plain}
\newtheorem{theorem}{Theorem}[section]
\newtheorem{corollary}[theorem]{Corollary}
\newtheorem{lemma}[theorem]{Lemma}
\newtheorem{proposition}[theorem]{Proposition}
\newtheorem{qnbank}[theorem]{Question Bank}
\let\oldqnbank\qnbank
\renewcommand{\qnbank}{\oldqnbank\normalfont}
\newtheorem{process}[theorem]{Process}
\let\oldprocess\process
\renewcommand{\process}{\oldprocess\normalfont}
\newtheorem{method}[theorem]{Method}
\let\oldmethod\method
\renewcommand{\method}{\oldmethod\normalfont}
\newtheorem{definition}[theorem]{Definition}
\let\olddefinition\definition
\renewcommand{\definition}{\olddefinition\normalfont}
\newtheorem{example}[theorem]{Example}
\let\oldexample\example
\renewcommand{\example}{\oldexample\normalfont}
\newtheorem{remark}[theorem]{Remark}
\let\oldremark\remark
\renewcommand{\remark}{\oldremark\normalfont}
\newtheorem{notation}[theorem]{Notation}
\let\oldnotation\notation
\renewcommand{\notation}{\oldnotation\normalfont}
\newtheorem{principle}[theorem]{Principle}
\let\oldprinciple\principle
\renewcommand{\principle}{\oldprinciple\normalfont}
\newtheorem{mechanism}[theorem]{Mechanism}
\let\oldmechanism\mechanism
\renewcommand{\mechanism}{\oldmechanism\normalfont}
\newtheorem{taxonomy}[theorem]{Taxonomy}
\let\oldtaxonomy\taxonomy
\renewcommand{\taxonomy}{\oldtaxonomy\normalfont}
\newtheorem{model}[theorem]{Model}
\let\oldmodel\model
\renewcommand{\model}{\oldmodel\normalfont}
\newtheorem{assumption}[theorem]{Assumption}
\let\oldassumption\assumption
\renewcommand{\assumption}{\oldassumption\normalfont}
\newtheorem{casestudy}[theorem]{Case Study}
\let\oldcasestudy\casestudy
\renewcommand{\casestudy}{\oldcasestudy\normalfont}
\newtheorem{empiricalfinding}[theorem]{Empirical Finding}
\let\oldempiricalfinding\empiricalfinding
\renewcommand{\empiricalfinding}{\oldempiricalfinding\normalfont}
\newtheorem{practitionernote}[theorem]{Practitioner Note}
\let\oldpractitionernote\practitionernote
\renewcommand{\practitionernote}{\oldpractitionernote\normalfont}
\newtheorem{caveat}[theorem]{Caveat}
\let\oldcaveat\caveat
\renewcommand{\caveat}{\oldcaveat\normalfont}
\newtheorem{marketrule}[theorem]{Rule}
\let\oldmarketrule\marketrule
\renewcommand{\marketrule}{\oldmarketrule\normalfont}
\newtheorem{takeaway}[theorem]{Takeaway}
\let\oldtakeaway\takeaway
\renewcommand{\takeaway}{\oldtakeaway\normalfont}

%%%%%%%%%%%%%%%%%%%%%%%%%%%%%%%%%%%%%%%%%%%%%%%%%%%%%%%%%%%%%%%%%%%%%%%%%%%%%%%
% Breakable Algorithm Environment (Algorithms spanning multiple pages)
%%%%%%%%%%%%%%%%%%%%%%%%%%%%%%%%%%%%%%%%%%%%%%%%%%%%%%%%%%%%%%%%%%%%%%%%%%%%%%%

\makeatletter
\newenvironment{breakablealgorithm}
 {%
  \begin{center}
    \refstepcounter{algorithm} % Increment algorithm counter
    \hrule height.8pt depth0pt \kern2pt
    \renewcommand{\caption}[2][\relax]{%
      {\raggedright\textbf{Algorithm \thealgorithm:} ##2\par}%
      \ifx\relax##1\relax
        \addcontentsline{loa}{algorithm}{\protect\numberline{\thealgorithm}##2}%
      \else
        \addcontentsline{loa}{algorithm}{\protect\numberline{\thealgorithm}##1}%
      \fi
      \kern2pt\hrule\kern2pt
    }%
 }{%
    \kern2pt\hrule
  \end{center}
 }
\makeatother

%%%%%%%%%%%%%%%%%%%%%%%%%%%%%%%%%%%%%%%%%%%%%%%%%%%%%%%%%%%%%%%%%%%%%%%%%%%%%%%
% Code Listings Style
%%%%%%%%%%%%%%%%%%%%%%%%%%%%%%%%%%%%%%%%%%%%%%%%%%%%%%%%%%%%%%%%%%%%%%%%%%%%%%%

\lstdefinestyle{codestyle}{
  commentstyle=\color{codegreen},
  keywordstyle=\color{magenta},
  numberstyle=\tiny\color{codegray},
  stringstyle=\color{codepurple},
  basicstyle=\ttfamily\normalsize,
  breakatwhitespace=false,
  breaklines=true,
  captionpos=b,
  keepspaces=true,
  numbers=left,
  numbersep=10pt,
  frame=single,
  showspaces=false,
  showstringspaces=false,
  showtabs=false,
  tabsize=2
}
\lstset{style=codestyle}

% -----------------------------
% Hyperlinks (loaded last: loaded before float and algorithm it duplicated figure anchors)
% -----------------------------
\usepackage{hyperref}
\hypersetup{
  colorlinks,
  linkcolor={blue!80!black},
  citecolor={green!80!black},
  urlcolor={red!80!black}
}

\usepackage[T1]{fontenc}
\usepackage{amsmath,amssymb,booktabs,graphicx,hyperref}
\hypersetup{colorlinks}
\setlength{\emergencystretch}{3em}   % lets a line holding a long key break
\newcounter{qfentry}[chapter]
\renewcommand{\theqfentry}{\thechapter.\arabic{qfentry}}
% \qftype{type}{Name} declares the environment qf<type>; its title wraps like a paragraph.
\newcommand{\qftype}[2]{\newenvironment{qf#1}[1][]{\par\addvspace{\medskipamount}%
  \refstepcounter{qfentry}{\raggedright\noindent\bfseries #2 \theqfentry
  \if\relax\detokenize{##1}\relax\else\quad ##1\fi\par}\nopagebreak}%
  {\par\addvspace{\medskipamount}}}
\qftype{definition}{Definition}\qftype{convention}{Convention}\qftype{result}{Result}
\qftype{model}{Model}\qftype{method}{Method}\qftype{mechanism}{Mechanism}
\qftype{evidence}{Evidence}\qftype{data}{Data}\qftype{rule}{Rule}
\qftype{pitfall}{Pitfall}\qftype{example}{Example}
\newcommand{\qffield}[1]{\par\textbf{#1.}\enspace}   % starts one of an entry's fields
% \citep[locator]{key} links the key to its line in the sources if it has one,
% and otherwise to its entry in the pinned catalogue PDF.
\newcommand{\qfcatalogue}{../guide/institutional_quantitative_finance_reading_list.pdf}
\NewDocumentCommand{\citep}{o o m}{[\IfValueT{#2}{#1 }%
  \ifcsname r@src:#3\endcsname\hyperref[src:#3]{\nolinkurl{#3}}%
  \else\href{\qfcatalogue\##3}{\nolinkurl{#3}}\fi
  \IfValueTF{#2}{, #2}{\IfValueT{#1}{, #1}}]}
% \qfsource{key}{reference} records one source outside the catalogue, in front/sources.tex.
\newcommand{\qfsource}[2]{\par\noindent\phantomsection\label{src:#1}\nolinkurl{#1}\quad #2\par}
% \qfsections{key, ...} inputs <key>/section.tex for each listed key that has one, in order.
\ExplSyntaxOn
\NewDocumentCommand{\qfsections}{m}
  {\clist_map_inline:nn {#1} {\InputIfFileExists{##1/section}{}{}}}
\ExplSyntaxOff

\pagestyle{plain}
\begin{document}
\chapter*{Portfolio construction and risk: legacy portfolio construction models}

\section*{Portfolio construction models (portfolio construction / source review / optimal execution / high-frequency strategies reviews)}
% Kept verbatim when this file was migrated on 2026-10-15 (notes/source-review-records.md, portfolio construction review): the whole of
% sysinv/portconstmod.tex below its migrate-to line, the introductory paragraph and 3 blocks in order with their
% prose and 13 items (4, 3 and 6); only the subsection is unnumbered here. No other text, table or figure was in it.
% Destinations: introduction -> meth:rule-based-position-sizing, mod:mean-variance-portfolio-optimisation;
% Rule-Based Models -> meth:rule-based-position-sizing; Optimisers Models Parameters ->
% mod:mean-variance-portfolio-optimisation; Optimisation Techniques -> mod:mean-variance-portfolio-optimisation,
% meth:signal-portfolio-allocation, stubs (Black-Litterman, resampled efficiency); the data-mining item not adopted.
% Sources: Narang (s7-narang), Inside the Black Box, third edition, Wiley 2024, chap. 6 "Portfolio Construction
% Models", pp. 99-125 (printed page = PDF page - 20), exhibits by page image; Boyd et al. (wide-boyd-markowitz-seventy),
% authors' version of 5 January 2024.

% From compendium/legacy/portfolio-construction-risk/sysinv/portconstmod.tex (whole file, deleted 2026-10-15 once empty).
\subsection*{Portfolio Construction Models}

Comes in two major forms: rule-based, optimisers. Rule-based models are based on heuristics, can be exceedingly simple or rather complex, and derived from human experience (trial and error). Optimisers comprises of an objective function and uses algorithms to reach the end goal.

\begin{definition} \hlt{Rule-Based Models}
\begin{enumerate}[label=\roman*.]
\setlength{\itemsep}{0pt}
\item Equal Position Weighting: used if portfolio manager believes that if a position is good enough to own, no other information is needed in determining its size. Strength of signal is not used as input in weighting.\\
Model assumes that there is sufficient statistical strength and power to predict not only direction but also magnitude relative to other forecasts in the portfolio. Portfolio takes few large bets on 'best' forecast, many smaller bets on less dramatic forecasts; may take excess risk in an idiosyncratic event on a seemingly attractive position, resulting in adverse selection bias.
\item Equal Risk Weighting: adjust position sizes inversely to volatilities or a measure of risk. More volatile positions given smaller allocations, less volatile positions given larger allocations.\\
When unit of risk is equalised, it is almost always a backward-looking measurement such as volatility. If volatility changes with time, then model will be misled.
\item Alpha-Driven Weighting: position size based primarily on alpha model. Alpha signal determines size of position, but usually with size limits. Constraints used also includes limits on size of total bet on a group. May also have a function that relates the magnitude of forecast to size of position.\\
If model used in futures trend following, might suffer sharp drawdowns. Reliance on accuracy of alpha.
\item Decision-Tree Weighting: decision path to arrive at the allocation for given instrument, depending on type of alpha model and type of instrument. Constraints may include percentage limits for allocation.\\
Model size grows dramatically if more alpha models or mode types of positions are included.
\end{enumerate}
\end{definition}

\begin{remark} \hlt{Optimisers Models Parameters}\\
Harry Markowitz's mean variance optimisation (MVO) as the pioneer model. Models are based on principles of modern portfolio theory (MPT). Inputs include asset expected return (mean), asset variance, expected correlation matrix. Other inputs include size of portfolio in currency terms, desired risk level (volatility or expected drawdown), and other constraints such as liquidity, universe limits.\\
Model uses an objective function and an algorithm to seek the goal, usually maximising return of portfolio relative to volatility of portfolio returns.
\begin{enumerate}[label=\roman*.]
\setlength{\itemsep}{0pt}
\item Expected Return: alpha models as basis of expected return, which also includes expected direction.
\item Expected Volatility: stochastic volatility forecasting methods is commonly used, as volatility may have high and low periods, with occasional jumps. GARCH model is most used.
\item Expected Correlation: as instrument correlations are not stable over time, it is more appropriate to group assets together before computing correlation within the group.
\end{enumerate}
\end{remark}

\begin{method} \hlt{Optimisation Techniques}
\begin{enumerate}[label=\roman*.]
\setlength{\itemsep}{0pt}
\item Unconstrained Optimisation: most basic form with no constraints. Might provide a single-instrument portfolio, where all money will be invested in instrument with highest risk-adjusted return.
\item Constrained optimisation: constraints include position limits, limits on various groupings of instruments. Might result in constraints driving the portfolio construction more than the optimiser.
\item Black-Litterman Optimisation: blends investor expectations with a degree of confidence about those expectations, and these with historical precedent evident in the data. Adjusts historically observed correlation levels by utilising investor's forecast of return for the various instruments.
\item Grinold and Kahn's Approach: builds a portfolio of signals, instead of sizing positions. To build factor portfolios, each of which are usually rule-based portfolios based on a single type of alpha forecast. Each portfolio backtested, then series of returns are then treated as instruments of a portfolio by the optimiser.\\
Number of factor portfolios is more manageable, usually not more than 20. What is optimised is then a handful of factor portfolios. The model allows for inclusion of risk model, transaction cost model, portfolio size, and risk targets as inputs.
\item Resampled Efficiency: to improve the inputs to optimisation by addressing oversensitivity to estimation error. To resample data using Monte Carlo simulation to reduce estimation error in inputs to the optimiser.
\item Data-Mining Approaches: machine learning techniques such as supervised learning or genetic algorithms used, as MVO involves searching many possible portfolios to find the best.
\end{enumerate}
\end{method}

\end{document}
